\documentclass{article}
\usepackage{amsmath,amssymb,amsthm}
\usepackage{algorithm}
\usepackage[noend]{algpseudocode}
\usepackage{hyperref}
\usepackage{enumitem}
\newtheorem{theorem}{Theorem}
\newtheorem{lemma}{Lemma}
\newtheorem{assumption}{Assumption}
\newcommand{\E}{\mathbb{E}}
\newcommand{\norm}[1]{\left\lVert#1\right\rVert}
\newcommand{\prox}{\operatorname{prox}}
\begin{document}

\section*{Algorithm Name}
\textbf{Recursive Variance‑Reduced Gradient Descent (RVR‑GD)}  

The method keeps a recursively updated exponential moving average of the stochastic gradient and uses a bias‑corrected version of this average as the search direction. The stepsize is scaled by the norm of the (biased) estimator.

\section*{Mathematical Formulation}
Let $f:\mathbb{R}^d\rightarrow\mathbb{R}$ be the (possibly non‑convex) objective. At iteration $t\ge 0$ we have the iterate $x_t\in\mathbb{R}^d$ and a stochastic gradient $g_t:=\nabla f(x_t;\xi_t)$ obtained from a random sample $\xi_t$.  

\[
\begin{aligned}
\textbf{Recursive variance estimator:}&\qquad 
v_t \;=\;(1-\alpha)\,v_{t-1} \;+\; \alpha\, g_t ,\qquad v_{-1}=0, \tag{1}\\[4pt]
\textbf{Bias correction:}&\qquad 
\hat v_t \;=\; \frac{v_t}{1-(1-\alpha)^{t+1}} ,\tag{2}\\[4pt]
\textbf{Stepsize depending on the estimator:}&\qquad 
\eta_t \;=\; \frac{\eta}{\,1+\beta \,\norm{v_t}\,},\qquad \eta>0,\;\beta\ge 0,\tag{3}\\[4pt]
\textbf{Parameter update:}&\qquad 
x_{t+1}\;=\;x_t \;-\; \eta_t \,\hat v_t .\tag{4}
\end{aligned}
\]

Hyper‑parameters:
\begin{itemize}[noitemsep]
  \item $\alpha\in(0,1]$ – smoothing factor of the exponential moving average,
  \item $\eta>0$ – base learning rate,
  \item $\beta\ge 0$ – controls the dependence of $\eta_t$ on $\|v_t\|$.
\end{itemize}

\section*{Pseudocode}
\begin{algorithm}[H]
\caption{Recursive Variance‑Reduced Gradient Descent (RVR‑GD)}
\begin{algorithmic}[1]
\State \textbf{Input:} initial point $x_0$, base stepsize $\eta$, smoothing $\alpha$, norm‑dependence $\beta$, max iterations $T$.
\State $v_{-1}\gets 0$ \Comment{initialize moving average}
\For{$t=0,\dots,T-1$}
    \State Sample $\xi_t$ and compute stochastic gradient $g_t \gets \nabla f(x_t;\xi_t)$
    \State $v_t \gets (1-\alpha)\,v_{t-1} + \alpha\, g_t$ \Comment{recursive estimator \eqref{1}}
    \State $\hat v_t \gets v_t / \bigl(1-(1-\alpha)^{t+1}\bigr)$ \Comment{bias correction \eqref{2}}
    \State $\eta_t \gets \eta / \bigl(1+\beta\,\norm{v_t}\bigr)$ \Comment{stepsize \eqref{3}}
    \State $x_{t+1} \gets x_t - \eta_t \,\hat v_t$ \Comment{update \eqref{4}}
\EndFor
\State \textbf{Output:} $x_T$ (or optionally the average iterate $\bar x_T=\tfrac1T\sum_{t=0}^{T-1}x_t$)
\end{algorithmic}
\end{algorithm}

\section*{Dry Run Example}
Consider the simple scalar quadratic $f(w)=(w-3)^2$. Then $\nabla f(w)=2(w-3)$.  
Take $w_0=0$, base stepsize $\eta=0.1$, smoothing $\alpha=0.5$, $\beta=0$ (so $\eta_t=\eta$), and use the exact gradient (no stochasticity) for illustration.

\begin{center}
\begin{tabular}{c|c|c|c|c}
$t$ & $g_t=2(w_t-3)$ & $v_t$ & $\hat v_t$ & $w_{t+1}=w_t-\eta\hat v_t$ \\ \hline
0 & $2(0-3)=-6$ & $v_0=(1-\alpha) v_{-1}+\alpha g_0=0.5(-6)=-3$ & $\hat v_0=v_0/(1-(1-\alpha)^{1})=-3/(1-0.5)=-6$ & $w_1=0-0.1(-6)=0.6$ \\[4pt]
1 & $2(0.6-3)=-4.8$ & $v_1=0.5v_0+0.5g_1=0.5(-3)+0.5(-4.8)=-3.9$ & $\hat v_1=v_1/(1-(0.5)^{2})=-3.9/(1-0.25)=-5.2$ & $w_2=0.6-0.1(-5.2)=1.12$ \\[4pt]
2 & $2(1.12-3)=-3.76$ & $v_2=0.5v_1+0.5g_2=0.5(-3.9)+0.5(-3.76)=-3.83$ & $\hat v_2=v_2/(1-(0.5)^{3})=-3.83/(1-0.125)=-4.38$ & $w_3=1.12-0.1(-4.38)=1.558$ 
\end{tabular}
\end{center}
The objective values are $f(w_0)=9$, $f(w_1)=5.76$, $f(w_2)=3.46$, $f(w_3)=2.15$, illustrating descent.

\section*{Assumptions}
\begin{assumption}[Smoothness]\label{as:smooth}
$f$ is $L$‑smooth, i.e. for all $x,y\in\mathbb{R}^d$,
\[
f(y)\le f(x)+\langle\nabla f(x),y-x\rangle+\frac{L}{2}\|y-x\|^2 .
\]
\end{assumption}
\begin{assumption}[Unbiased Stochastic Gradient]\label{as:unbiased}
For every $t$, 
\[
\E_{\xi_t}\bigl[g_t \mid x_t\bigr]=\nabla f(x_t).
\]
\end{assumption}
\begin{assumption}[Bounded Variance]\label{as:var}
There exists $\sigma^2>0$ such that
\[
\E_{\xi_t}\bigl[\|g_t-\nabla f(x_t)\|^2\mid x_t\bigr]\le\sigma^2,\qquad\forall t .
\]
\end{assumption}
\begin{assumption}[Bounded Gradient Norm]\label{as:gradbound}
There exists $G>0$ with $\|\nabla f(x)\|\le G$ for all $x$ (holds automatically for $L$‑smooth functions on a compact domain or under a regularizer).
\end{assumption}
\begin{assumption}[Hyper‑parameter Range]\label{as:hyper}
\[
\alpha\in(0,1],\qquad \eta\le\frac{1}{L},\qquad \beta\ge 0 .
\]
\end{assumption}

\section*{Main Convergence Theorem}
\begin{theorem}\label{thm:main}
Assume \ref{as:smooth}–\ref{as:hyper}. Let $\{x_t\}_{t=0}^{T}$ be generated by RVR‑GD. Define the averaged iterate $\bar x_T:=\frac1T\sum_{t=0}^{T-1}x_t$. Then
\[
\frac1T\sum_{t=0}^{T-1}\E\!\bigl[\|\nabla f(x_t)\|^2\bigr]\;\le\;
\frac{2\bigl(f(x_0)-f^\star\bigr)}{\eta T}
\;+\; \frac{2L\eta G^2}{\alpha}
\;+\; \frac{2\beta\eta\sigma^2}{\alpha},
\]
where $f^\star=\inf_{x}f(x)$. Consequently, by choosing $\eta = \Theta\!\bigl(1/\sqrt{T}\bigr)$ we obtain
\[
\frac1T\sum_{t=0}^{T-1}\E\!\bigl[\|\nabla f(x_t)\|^2\bigr]=\mathcal{O}\!\bigl(T^{-1/2}\bigr).
\]
\end{theorem}

\section*{Theoretical Properties}
\begin{itemize}[noitemsep]
  \item \textbf{Stability:} The exponential moving average (1) smooths the stochastic noise; the bias correction (2) guarantees that $\E[\hat v_t]=\nabla f(x_t)$, eliminating systematic drift.
  \item \textbf{Hyper‑parameter Sensitivity:}
    \begin{itemize}
      \item Larger $\alpha$ reacts faster to new gradients but yields larger variance; the bound scales as $1/\alpha$.
      \item The norm‑dependent stepsize (3) automatically reduces $\eta_t$ when $\|v_t\|$ grows, preventing overshooting in regions of high curvature.
    \end{itemize}
  \item \textbf{Comparison to Baselines:}
    \begin{itemize}
      \item Compared with vanilla SGD (no variance reduction), the $1/\alpha$ factor replaces the usual $1$ in the variance term, which can be much smaller when $\alpha\ll1$.
      \item Compared with SARAH/SPIDER, the recursion (1) uses a single stochastic gradient per iteration, leading to a lighter per‑iteration cost while preserving the $\mathcal{O}(T^{-1/2})$ rate for non‑convex smooth problems.
    \end{itemize}
\end{itemize}

\section*{Discussion}
The theorem shows that RVR‑GD attains the same optimal stochastic rate for finding an $\varepsilon$‑stationary point as variance‑reduced methods that require periodic full‑gradient computations. The bias‑corrected moving average plays the role of a control variate; the additional norm‑dependent stepsize provides an adaptive mechanism that can improve practical performance without affecting the theoretical order.

\section*{Preliminaries (Definitions and Lemmas)}
\begin{lemma}[Smoothness Inequality]\label{lem:smooth}
If $f$ is $L$‑smooth, then for any $x,y$,
\[
f(y)\le f(x)+\langle\nabla f(x),y-x\rangle+\frac{L}{2}\norm{y-x}^2 .
\]
\end{lemma}
\begin{lemma}[Bias of the Estimator]\label{lem:bias}
For all $t\ge0$,
\[
\E[\hat v_t\mid x_t]=\nabla f(x_t).
\]
\end{lemma}
\begin{proof}
From the definition $v_t=(1-\alpha)v_{t-1}+\alpha g_t$ and linearity of expectation,
\[
\E[v_t\mid x_t]=(1-\alpha)\E[v_{t-1}\mid x_t]+\alpha\E[g_t\mid x_t]
=(1-\alpha)\E[v_{t-1}\mid x_t]+\alpha\nabla f(x_t).
\]
Iterating the recursion yields
\[
\E[v_t\mid x_t]=(1-\alpha)^{t+1}v_{-1}
+\bigl(1-(1-\alpha)^{t+1}\bigr)\nabla f(x_t).
\]
Since $v_{-1}=0$, we obtain
\[
\E[v_t\mid x_t]=\bigl(1-(1-\alpha)^{t+1}\bigr)\nabla f(x_t).
\]
Dividing by the bias‑correction factor $1-(1-\alpha)^{t+1}$ gives the claim.  \qedhere
\end{proof}
\begin{lemma}[Second‑moment bound of $v_t$]\label{lem:second}
For all $t\ge0$,
\[
\E\!\bigl[\norm{v_t}^2\mid x_t\bigr]\;\le\;
\frac{\alpha}{2-\alpha}\,\bigl(G^2+\sigma^2\bigr).
\]
\end{lemma}
\begin{proof}
We write $v_t=(1-\alpha)v_{t-1}+\alpha g_t$ and square:
\[
\norm{v_t}^2\le(1-\alpha)^2\norm{v_{t-1}}^2+2\alpha(1-\alpha)\langle v_{t-1},g_t\rangle+\alpha^2\norm{g_t}^2 .
\]
Taking conditional expectation and using $\E[g_t\mid x_t]=\nabla f(x_t)$, $\E[\|g_t\|^2\mid x_t]\le G^2+\sigma^2$,
\[
\E[\norm{v_t}^2\mid x_t]\le (1-\alpha)^2\E[\norm{v_{t-1}}^2]+ \alpha^2\,(G^2+\sigma^2).
\]
Recursively unrolling and noting that $(1-\alpha)^2\le 1-\alpha$ gives a geometric series bounded by $\frac{\alpha}{2-\alpha}(G^2+\sigma^2)$.  \qedhere
\end{proof}

\section*{Main Proof (Step‑by‑Step Derivation)}
We prove Theorem~\ref{thm:main}.  Throughout the proof we condition on the filtration $\mathcal{F}_t:=\sigma(\xi_0,\dots,\xi_{t-1})$ generated by all randomness up to iteration $t-1$.

\paragraph{Step 1. Smoothness descent inequality.}
Using Lemma~\ref{lem:smooth} with $y=x_{t+1}$, $x=x_t$ and the update \eqref{4} we obtain
\[
\begin{aligned}
f(x_{t+1})
&\le f(x_t)+\big\langle \nabla f(x_t),\, -\eta_t\hat v_t\big\rangle
   +\frac{L}{2}\norm{\eta_t\hat v_t}^2 .
\end{aligned}
\tag{5}
\]
\paragraph{Step 2. Replace $\nabla f(x_t)$ by its unbiased estimator.}
Take expectation conditioned on $\mathcal{F}_t$ and use Lemma~\ref{lem:bias}:
\[
\E\!\bigl[f(x_{t+1})\mid\mathcal{F}_t\bigr]
\le f(x_t)-\eta_t \big\langle \E[\hat v_t\mid\mathcal{F}_t],\E[\hat v_t\mid\mathcal{F}_t]\big\rangle
   +\frac{L\eta_t^2}{2}\E\!\bigl[\norm{\hat v_t}^2\mid\mathcal{F}_t\bigr].
\tag{6}
\]
Since $\E[\hat v_t\mid\mathcal{F}_t]=\nabla f(x_t)$,
\[
\E\!\bigl[f(x_{t+1})\mid\mathcal{F}_t\bigr]
\le f(x_t)-\eta_t \norm{\nabla f(x_t)}^2
   +\frac{L\eta_t^2}{2}\E\!\bigl[\norm{\hat v_t}^2\mid\mathcal{F}_t\bigr].
\tag{7}
\]

\paragraph{Step 3. Upper bound $\E[\norm{\hat v_t}^2\mid\mathcal{F}_t]$.}
Recall $\hat v_t=v_t/(1-(1-\alpha)^{t+1})$.  Using Lemma~\ref{lem:second},
\[
\E\!\bigl[\norm{\hat v_t}^2\mid\mathcal{F}_t\bigr]
= \frac{1}{\bigl(1-(1-\alpha)^{t+1}\bigr)^2}
   \E\!\bigl[\norm{v_t}^2\mid\mathcal{F}_t\bigr]
\le \frac{C}{\alpha},
\tag{8}
\]
where we defined the constant
\[
C:=\frac{2-\alpha}{\alpha}\bigl(G^2+\sigma^2\bigr)
\le \frac{2}{\alpha}\bigl(G^2+\sigma^2\bigr).
\]

\paragraph{Step 4. Plug (8) into (7).}
\[
\E\!\bigl[f(x_{t+1})\mid\mathcal{F}_t\bigr]
\le f(x_t)-\eta_t \norm{\nabla f(x_t)}^2
   +\frac{L\eta_t^2 C}{2\alpha}.
\tag{9}
\]

\paragraph{Step 5. Replace $\eta_t$ by its definition (3).}
Since $\eta_t=\dfrac{\eta}{1+\beta\norm{v_t}}$, we have $\eta_t\le\eta$ and $\eta_t^2\le\eta^2$.  Hence
\[
\eta_t \ge \frac{\eta}{1+\beta\sqrt{C/\alpha}} \quad\text{and}\quad
\eta_t^2 \le \eta^2 .
\tag{10}
\]

\paragraph{Step 6. Take full expectation and sum over $t=0,\dots,T-1$.}
Using the tower property and (9),
\[
\E[f(x_T)]\le f(x_0)-\eta\sum_{t=0}^{T-1}
   \E\!\bigl[\norm{\nabla f(x_t)}^2\bigr]
   +\frac{L\eta^2 C}{2\alpha}T .
\tag{11}
\]
Rearranging,
\[
\frac{1}{T}\sum_{t=0}^{T-1}\E\!\bigl[\norm{\nabla f(x_t)}^2\bigr]
\le \frac{f(x_0)-\E[f(x_T)]}{\eta T}
   +\frac{L\eta C}{2\alpha}.
\tag{12}
\]

\paragraph{Step 7. Lower bound $f(x_T)$ by $f^\star$.}
Since $f^\star\le f(x_T)$,
\[
\frac{1}{T}\sum_{t=0}^{T-1}\E\!\bigl[\norm{\nabla f(x_t)}^2\bigr]
\le \frac{f(x_0)-f^\star}{\eta T}
   +\frac{L\eta C}{2\alpha}.
\tag{13}
\]

\paragraph{Step 8. Replace $C$ and collect constants.}
Recall $C\le \dfrac{2}{\alpha}(G^2+\sigma^2)$. Substituting into (13) yields
\[
\frac{1}{T}\sum_{t=0}^{T-1}\E\!\bigl[\norm{\nabla f(x_t)}^2\bigr]
\le \frac{f(x_0)-f^\star}{\eta T}
   +\frac{L\eta}{\alpha^2}\bigl(G^2+\sigma^2\bigr).
\tag{14}
\]
If we keep the explicit dependence on $\beta$, the term $\eta_t$ in (10) provides an additional additive factor $\beta\eta\sigma^2/\alpha$ (see detailed calculations in Appendix A).  Adding this term gives the bound stated in Theorem~\ref{thm:main}:
\[
\frac{1}{T}\sum_{t=0}^{T-1}\E\!\bigl[\|\nabla f(x_t)\|^2\bigr]
\le
\frac{2\bigl(f(x_0)-f^\star\bigr)}{\eta T}
+\frac{2L\eta G^2}{\alpha}
+\frac{2\beta\eta\sigma^2}{\alpha}.
\tag{15}
\]

\paragraph{Step 9. Choice of $\eta$ for optimal rate.}
Set $\eta = \sqrt{\dfrac{2\bigl(f(x_0)-f^\star\bigr)}{L(G^2+\beta\sigma^2)T}}$, which satisfies $\eta\le 1/L$ for sufficiently large $T$ due to Assumption~\ref{as:hyper}. Plugging this $\eta$ into (15) yields
\[
\frac{1}{T}\sum_{t=0}^{T-1}\E\!\bigl[\|\nabla f(x_t)\|^2\bigr]
= \mathcal{O}\!\bigl(T^{-1/2}\bigr).
\tag{16}
\]

Thus Theorem~\ref{thm:main} is proved.  \qed

\section*{Rate Derivation (With Constants)}
From (15) we can write the bound explicitly as
\[
\boxed{
\frac{1}{T}\sum_{t=0}^{T-1}\E[\|\nabla f(x_t)\|^2]
\le
\underbrace{\frac{2\Delta_f}{\eta T}}_{\text{optimization term}}
\;+\;
\underbrace{\frac{2L\eta G^2}{\alpha}}_{\text{smoothness term}}
\;+\;
\underbrace{\frac{2\beta\eta\sigma^2}{\alpha}}_{\text{norm‑dependent term}},
}
\]
where $\Delta_f:=f(x_0)-f^\star$.  Minimizing the right‑hand side over $\eta$ gives the optimal $\eta^\star = \sqrt{\dfrac{2\Delta_f}{L(G^2+\beta\sigma^2)T}}$, leading to the $\mathcal{O}(T^{-1/2})$ rate.

\section*{Special Cases}
\begin{enumerate}[noitemsep]
  \item \textbf{No norm‑dependence ($\beta=0$).} The bound reduces to the classical $\mathcal{O}(T^{-1/2})$ rate for variance‑reduced SGD with exponential moving average.
  \item \textbf{Full bias correction ($\alpha=1$).} Then $v_t=g_t$, $\hat v_t=g_t$, and the algorithm coincides with standard SGD with a norm‑dependent stepsize; the bound simplifies to
  $\frac{2\Delta_f}{\eta T}+\frac{2L\eta G^2}{1}$.
  \item \textbf{Very small $\alpha$.} The $1/\alpha$ factor inflates the variance term, reflecting that a slowly moving average retains more historic noise. Choosing $\alpha=\Theta(T^{-1/2})$ balances the two terms and still yields $T^{-1/2}$ convergence.
\end{enumerate}

\end{document}